\section{同一个模型可以既公平又不公平}
\label{sec:17-Constrained-Guarantees/03-Fairness-Impossibility/01}

一个再犯风险评分系统上线一年后，两份分析报告同时摆上桌，结论相反。

开发方那份是这样算的：把所有被评分的人按分数档分组，在每一档内部分别统计两个群体两年内的实际再犯比例。两条曲线几乎重合——分数为 \(7\) 的人，不论属于哪个群体，再犯比例都在 \(0.6\) 附近。分数这个数字对两边含义一致，没有一边被系统性地抬高。

批评方那份换了一个切法：只取两年内确实没有再犯的人，问他们当初被判为高风险的比例。一个群体是 \(45\%\)，另一个是 \(23\%\)。在同样没有再犯的人里，一边被判高风险的机会接近另一边的两倍。

两份报告没有一份算错，它们甚至读的是同一张表。分歧不在数字上，在于\textbf{条件写在竖线的哪一边}。

\subsection{同一张表，两个方向的条件概率}

记 \(A\) 为群体属性，\(Y\in\{0,1\}\) 为两年内是否再犯这个真实结果，\(S\) 为模型输出的分数，\(\hat Y=\ind\{S\ge t\}\) 为按阈值 \(t\) 得到的二元判定。

开发方检验的是
\[
P(Y=1\given S=s,\ A=a)\quad\text{与 }a\text{ 无关，对每个 }s,
\]
即给定分数去问真实结果的分布，这叫做\textbf{校准}。批评方检验的是
\[
P(\hat Y=1\given Y=0,\ A=a)\quad\text{与 }a\text{ 无关},
\]
即给定真实结果去问判定的分布，这是\textbf{假阳率}相等。

两个式子的条件量与被条件量恰好对调。它们描述的是联合分布 \(P(Y,\hat Y\given A)\) 的两种不同边缘化方式，彼此之间没有任何逻辑蕴含关系——一个成立不推出另一个成立，也不推出另一个不成立。所以``模型是校准的''与``模型的假阳率不等''可以同时为真，而且在现实数据上通常就是同时为真。

这种方向上的混淆本书已经遇到过一次。\hyperref[sec:11-Mathematical-Statistics/05-Hypothesis-Testing/02]{``\(p\) 值不是原假设为真的概率''}讲的是 \(P(\text{数据}\given H_0)\) 与 \(P(H_0\given\text{数据})\) 不是一回事，把它们互换会让人对一个从未计算过的量充满信心。Bayes 公式给出的正是两个方向之间那道换算关系，而换算需要一个基率——\hyperref[sec:05-Probability/02-Conditional-Probability/02]{``全概率公式与 Bayes 更新''}里那个医学检验的例子，说的就是同一个灵敏度在不同患病率下会导出完全不同的阳性预测值。公平性争论里的``基率''，就是两个群体各自的 \(P(Y=1\given A=a)\)。这一点在下一节会成为整件事的支点。

\subsection{四个判据，四种朴素的正义直觉}

争论之所以持久，是因为参与争论的人手里各有一个判据，而每个判据都对应一句听上去无可辩驳的话。把它们并排写出来，就能看清它们在对同一张表的不同部分做归一化。

\textbf{人口平等}要求
\[
P(\hat Y=1\given A=a)\quad\text{与 }a\text{ 无关}.
\]
它背后的直觉是``两个群体被判为高风险的比例应当一样''。这个判据完全不看 \(Y\)，因此它甚至不需要知道谁真的再犯了；它的强项也在这里——不依赖标签，所以不会被有偏的标签污染。它的问题同样在这里：若两个群体的真实 \(P(Y=1\given A=a)\) 确实不同，强行拉平判定比例就必须在某一边放松、在另一边收紧。

\textbf{机会平等}要求真阳率相等，
\[
P(\hat Y=1\given Y=1,\ A=a)\quad\text{与 }a\text{ 无关}.
\]
直觉是``真的会再犯的人，被识别出来的机会不该因群体而异''。它只约束 \(Y=1\) 那一行，对 \(Y=0\) 的人怎么被对待不置一词。

\textbf{等化几率}把两行一起管住：真阳率与假阳率都与 \(a\) 无关。直觉是``判定只应当依赖真实结果，不应当在真实结果之外还携带群体信息''。用条件独立的语言写就是 \(\hat Y\perp A\given Y\)，这正是\hyperref[sec:05-Probability/02-Conditional-Probability/04]{``条件独立与信息结构''}里那个结构：给定 \(Y\) 之后，\(A\) 不再对 \(\hat Y\) 提供任何额外信息。

\textbf{校准}要求 \(P(Y=1\given S=s,A=a)\) 与 \(a\) 无关。直觉是``分数这个数字对所有人意思一样''。它管的是列——固定分数档，看这一档里 \(Y\) 怎么分布。校准是概率输出可用的前提：\hyperref[sec:11-Mathematical-Statistics/07-Statistical-Decision-and-Reliable-Prediction/04]{``Proper scoring rule 与校准衡量概率是否诚实''}说明了，没有校准的分数拿去做期望损失计算会给出错的行动，而这恰恰是评分系统被使用的方式。

\begin{center}
\begin{tikzpicture}[x=1cm,y=1cm,font=\small,>=stealth]
  % 2x2 表
  \draw[black!50,fill=blue!7]  (0,1.9) rectangle (1.9,3.8);
  \draw[black!50,fill=blue!14] (1.9,1.9) rectangle (3.8,3.8);
  \draw[black!50,fill=blue!14] (0,0) rectangle (1.9,1.9);
  \draw[black!50,fill=blue!7]  (1.9,0) rectangle (3.8,1.9);
  \node at (0.95,2.85) {漏判};
  \node at (2.85,2.85) {命中};
  \node at (0.95,0.95) {正确放行};
  \node at (2.85,0.95) {误判};
  % 表头
  \node[font=\footnotesize] at (0.95,4.15) {\(\hat Y=0\)};
  \node[font=\footnotesize] at (2.85,4.15) {\(\hat Y=1\)};
  \node[font=\footnotesize,anchor=east] at (-0.12,2.85) {\(Y=1\)};
  \node[font=\footnotesize,anchor=east] at (-0.12,0.95) {\(Y=0\)};
  % 行方向箭头
  \draw[<->,black!60] (0.35,3.5) -- (3.45,3.5);
  \draw[<->,black!60] (0.35,0.4) -- (3.45,0.4);
  % 列方向箭头
  \draw[<->,black!60] (3.45,3.15) -- (3.45,1.15);
  % 注释
  \node[anchor=west,align=left,font=\footnotesize,black!75] at (4.2,3.5)
    {\(\div\) 这一行的总数\\机会平等：真阳率相等};
  \node[anchor=west,align=left,font=\footnotesize,black!75] at (4.2,0.4)
    {\(\div\) 这一行的总数\\等化几率再加上：假阳率相等};
  \node[anchor=west,align=left,font=\footnotesize,black!75] at (4.2,2.0)
    {\(\div\) 这一列的总数\\校准：每个分数档内 \(Y\) 的比例相等};
  \node[anchor=west,align=left,font=\footnotesize,black!75] at (4.2,-0.75)
    {\(\div\) 全表总数\\人口平等：判为正例的比例相等};
  \draw[<->,black!60] (0.35,-0.35) -- (3.45,-0.35);
\end{tikzpicture}

\emph{四个判据读的是同一张 \(2\times2\) 表，区别只在于除以哪个分母——行、列，还是全表。分母不同，两个群体``相等''的含义就不同。}
\end{center}

这张图要记住的是分母。争论双方各自拿着一个真数字，却在心里除着不同的东西，于是谁也说服不了谁——不是因为谁不讲理，是因为他们在比较两个不同的量。开发方那份报告除的是列，批评方那份除的是行。

\subsection{为什么不能``都要''}

看到这里，很自然的反应是把四个判据都写进验收标准。这个反应值得记下来，因为下一节要证明的正是它不可行：只要两个群体的基率不同，且分类器不完美，校准与等化几率就不能同时成立。这不是当前模型不够好，也不是数据不够多——它是这张 \(2\times2\) 表的代数结构本身给出的限制。

在此之前先把一件事说清楚，它影响后面所有讨论怎么读。上面四个判据全都只依赖联合分布 \(P(A,S,Y)\)，都是可以从观察数据里算出来的量。一个判据被满足或被违反，是一句关于分布的陈述，不是一句关于机制的陈述。分布相同而机制不同的两种世界，在这些判据眼里完全一样——这一点在本单元第三节会成为主要话题。

还有一个技术性的但很常被忽略的前提：\(Y\) 是``两年内是否再犯''这个被记录下来的变量，不是``是否再犯''本身。逮捕记录经过了执法密度、报案率、司法流程的层层过滤。任何以 \(Y\) 为准绳的判据，检验的都是模型对这份记录的一致性，而不是对被记录之事的一致性。这与\hyperref[sec:13-Machine-Learning/01-Learning-Problems-and-Evaluation/02]{``损失、风险、目标与指标不是同一个量''}讲的是同一件事：指标是目标的代理，代理与目标之间的缝隙不会因为指标算得精确而缩小。

下一节把``不能都要''从一句观察变成一条命题，并把证明写完整。它只需要三四行代数，但结论管得很宽：它说明这场争论无法靠改进模型来终结，只能靠先说清要哪一个来推进。
